\uuid{77ZB}
\titre{Dérivée en un point}
\auteur{bodin ; hanani ; mzari}
\organisation{exo7}
\licence{CC BY-NC-SA 4.0}
\datecreate{2020-01-20}
\niveau{L1}
\module{Analyse}
\chapitre{Dérivabilité des fonctions réelles}
\sousChapitre{Autre}
\difficulte{2}
\theme{Dérivées, dérivée en un point}
\exo7qcmid{lille-1.derivees.015}

\contenu{
\begin{question}
\qclassification{124.00}
\qtags{motcle=dérivée en un point}
Soit $f(x)=\left\{\begin{array}{cl}\displaystyle x+x^2\sin \frac{1}{x}&\mbox{si }x\neq 0\\ \\ 0&\mbox{si }x=0.\end{array}\right.$

Quelles sont les bonnes réponses ?
\begin{answers}
    \bad{$f$ n'est pas dérivable en $0$.}
    \bad{$f$ est dérivable en $0$ est $f'(0)=0$.}
    \good{$f$ est dérivable en $0$ est $f'(0)=1$.}
    \good{Pour $x\neq 0$, $\displaystyle f'(x)=1+2x\sin \frac{1}{x}-\cos \frac{1}{x}$.}
\end{answers}
\begin{explanations}
On a
$$\lim_{x\to 0}\frac{f(x)-f(0)}{x-0}=\lim_{x\to 0}\left(1+x\sin \frac{1}{x}\right)=1.$$
Donc, $f$ est dérivable en $0$ et $f'(0)=1$. Par ailleurs, les règles de calcul donnent, pour $x\neq 0$,
$$f'(x)=(x)'+(x^2)'\sin \frac{1}{x}+x^2\left(\sin \frac{1}{x}\right)'=1+2x\sin \frac{1}{x}-\cos \frac{1}{x}.$$
\end{explanations}
\end{question}
}
